\subsection{Ambient families and vertical transversality}

An ambient family is the correct starting object; a tube is a conclusion about a
particular incidence inside it.

\begin{definition}[Ambient AEG family]
\label{def:ambient-aeg-family}
An \emph{ambient regular AEG family} over a smooth manifold $B$ consists of a
smooth fiber bundle $q:X\to B$ whose vertical tangent bundle $T^VX$ carries a
smooth orientation, a smooth vertical Riemannian metric $g^V$, and a smooth
function $A:X\to\R$ such that
$(M_b,g_b,A_b;\mu,\lambda)$ is a regular AES for every $b\in B$.  Its assignment
incidence is $\calZ(A)=A^{-1}(0)$.

Equivalently, the structure group of the surface bundle is reduced to
orientation-preserving diffeomorphisms.  Merely choosing an orientation on each
fiber without requiring those choices to vary continuously would not orient
$T^VX$ and is not sufficient here.

More generally, a \emph{field family over an AEG background} is a rank-$r$ real
vector bundle $E\to X$ with a smooth section $s$.  It need not be the real
assignment and need not satisfy the eikonal equation.  A complex-valued field is
the case $r=2$ after forgetting complex scalar multiplication.
\end{definition}

Let $T^VX=\ker(\dd q)$ be the vertical tangent bundle.  At a zero of $s$, choose
any connection to split $TX$ and define the vertical derivative
\[
 D^Vs:T^V_xX\longrightarrow E_x.
\]
Its value at a zero is independent of the chosen local trivialization or
connection: it is the derivative of the restricted section on the fiber.

\begin{theorem}[Parameterized zero-section theorem]
\label{thm:parameterized-zero-section}
Let $q:X\to B$ be a smooth fiber bundle and $E\to X$ a rank-$r$ vector bundle.
Suppose the section $s\in\Gamma(E)$ is vertically transverse to the zero section,
meaning that $D^Vs$ is surjective at every $x\in\calZ(s)$.  Then:
\begin{enumerate}
\item $\calZ(s)$ is a smooth submanifold of $X$ of codimension $r$;
\item the restricted projection
\[
 \pi=q|_{\calZ(s)}:\calZ(s)\longrightarrow B
 \]
is a smooth submersion;
\item if $\pi$ is surjective and proper, then it is a locally trivial smooth
fiber bundle.
\end{enumerate}
The third conclusion has the standard relative form when $X$ has boundary and the
zero incidence is neat, with the boundary restriction also a submersion.
\end{theorem}

\begin{proof}
Vertical surjectivity implies surjectivity of the full derivative of $s$, so the
zero-section transversality theorem gives the first conclusion.  To prove the
second, take $\xi\in T_bB$ and a lift $w\in T_xX$.  Since $D^Vs$ is surjective,
there is $v\in T_x^VX$ with $D^Vs(v)=-Ds(w)$.  Then
$w+v\in T_x\calZ(s)$ and $Dq(w+v)=\xi$, so $D\pi$ is surjective.  The third
conclusion is Ehresmann's fibration theorem; the neat relative version follows by
the corresponding boundary argument \cite{Ehresmann1952Connexions}.
\end{proof}

The rank is decisive.  If the fibers of $q$ have dimension $m$, then a regular
zero fiber has dimension $m-r$.

\begin{center}
\small
\begin{tabular}{cccc}
\toprule
Fiber dimension & Field rank & Fiberwise zero & Natural transport \\
\midrule
$2$ & $1$ & curves & component and mapping-class monodromy \\
$2$ & $2$ & points & permutation and braid monodromy \\
\bottomrule
\end{tabular}
\end{center}

For an ambient regular AEG family, equation~\eqref{eq:aeg-eikonal} makes
$D^VA=\dd A_b$ nonzero everywhere, so the first two conclusions are automatic at
zero.  Properness is not.

\subsection{The discriminant is not a single equation}

For finite-dimensional polynomial root families, the algebraic discriminant
detects repeated roots.  A general AEG family has more ways to lose topological
stability.

\begin{definition}[Structural discriminant]
\label{def:structural-discriminant}
For a parameterized zero problem, the \emph{structural discriminant} is the union
of the following loci in the parameter space, whenever they are present:
\begin{align*}
 \calD_{\mathrm{crit}}
   &=\{b:\text{$D^Vs$ fails to be surjective at some zero over $b$}\},\\
 \calD_{\partial}
   &=\{b:\text{the zero incidence fails the specified neat boundary condition}\},\\
 \calD_{\mathrm{sing}}
   &=\{b:\text{a zero meets the declared singular stratum or the metric degenerates}\},\\
 \calD_{\mathrm{dom}}
   &=\{b:\text{the domain, rank, or regularity class of the field changes}\},\\
 \calD_{\infty}
   &=\{b:\text{$\pi$ is not proper over any neighborhood of $b$}\}.
\end{align*}
If projective continuation is used, its ordinary poles and points at infinity are
recorded separately rather than silently absorbed into $\calD_{\mathrm{sing}}$.
\end{definition}

This definition is a bookkeeping union, not a claim that every stratum is a smooth
hypersurface.  The geometry of its intersections and a canonical Whitney
stratification remain open in general.

For the monic relative divisor of
Theorem~\ref{thm:relative-zero-divisor-decomposition}, the algebraic discriminant
is the scheme-theoretic analogue of the collision/ramification locus, and its
complement supports a finite \'etale root covering.  If $K\subset\C$, then on a
smooth analytic open of $B^{\mathrm{an}}$ where the polynomial is regarded as a
complex rank-$2$ field, this locus is the repeated-root, hence vertical-critical,
part of $\calD_{\mathrm{crit}}$.  Without that analytification and smooth-locus
restriction the two categories are not identified.  The algebraic discriminant
does not see $\calD_\partial$, $\calD_{\mathrm{sing}}$,
$\calD_{\mathrm{dom}}$, or $\calD_\infty$.  Thus arithmetic factorization and
geometric stability fit into the same framework without identifying the
polynomial discriminant with the entire structural discriminant.

\begin{corollary}[Stability off the discriminant]
\label{cor:stability-off-discriminant}
Let $U\subset B$ be connected.  Over $U$, assume that $q:X\to B$ and the
rank-$r$ bundle $E$ remain fixed smooth bundles, that $s$ is vertically transverse
to zero, and that the incidence is surjective and locally proper over $U$.  If a
boundary is present, also assume that the incidence is neat and that its boundary
projection is a submersion.  Then
$\pi:\calZ(s)|_U\to U$ is a smooth fiber bundle.  In particular, the diffeomorphism
type and number of connected components of the fiber are locally constant on $U$.
These are precisely the operative hypotheses encoded by saying that $U$ avoids
the relevant pieces of the structural discriminant.
\end{corollary}

\begin{proof}
For each $b\in U$, local properness supplies a neighborhood $V$ on which the
restricted incidence projection is proper.  The other hypotheses make that
restriction a surjective submersion, including the relative boundary condition
when needed.  Theorem~\ref{thm:parameterized-zero-section} therefore gives a
bundle chart over a possibly smaller neighborhood of $b$.  Such charts are exactly
the local-triviality assertion.
\end{proof}

\subsection{Escape to infinity without a critical zero}

The following regular AEG family shows why $\calD_\infty$ cannot be omitted.  Set
\[
 h_t(y)=y(ty-1),
 \qquad
 A(x,y,t)=e^x h_t(y),
 \]
and equip each copy of $\R^2$ with the smoothly varying conformal metric
\begin{equation}
 g_t=
 \frac{e^{2x}\bigl(h_t(y)^2+(2ty-1)^2\bigr)}
      {\mu^2+\lambda^2e^{2x}h_t(y)^2}
 (\dd x^2+\dd y^2).
 \label{eq:escape-family-metric}
\end{equation}
At a zero of $h_t$, its derivative $2ty-1$ equals $-1$ or $1$, so it never
vanishes.  More strongly, $h_t$ and $2ty-1$ never vanish simultaneously anywhere.
Thus every slice is a regular AES by Theorem~\ref{thm:conformal-realization}.

Its zero fibers are
\begin{equation}
 Z(A_t)=
 \begin{cases}
  \{y=0\},&t=0,\\
  \{y=0\}\amalg\{y=1/t\},&t\ne0.
 \end{cases}
 \label{eq:escape-zero-fibers}
\end{equation}
The total zero incidence is the disjoint union of the plane $y=0$ and the
hyperbolic sheet $ty=1$.  It is smooth, and its projection to $t$ is a submersion.
Nevertheless the second component escapes to spatial infinity as $t\to0$.

\begin{proposition}[Nonproper topology change]
\label{prop:nonproper-escape}
The family \eqref{eq:escape-family-metric} has no critical zero and no singular
metric, but the number of zero components changes at $t=0$.  Its incidence
projection is not proper over any neighborhood of $0$.
\end{proposition}

\begin{proof}
Only the last statement remains.  For $t_j\to0$ with $t_j\ne0$, the points
$(0,1/t_j,t_j)$ lie in the incidence above a compact parameter interval and have
no convergent subsequence.  Hence the inverse image of that interval is not
compact.
\end{proof}

Thus topology change does not require $\dd A_t=0$ at a finite zero.  It may occur at
infinity.  Conversely, a smooth total incidence does not imply a tube.

\subsection{Boundary control}

If a surface fiber has boundary, a zero curve may be closed or may end on the
boundary.  For a stable relative tube one requires the incidence to meet
$\partial^VX$ neatly and the restriction
$\pi:\partial\calZ\to B$ to remain a proper submersion.  Without this condition a
zero component can be created by tangency to the boundary even though the interior
field remains regular.  Such a parameter lies in $\calD_\partial$, not in
$\calD_{\mathrm{crit}}$.

\begin{problem}[Stratified AEG discriminant]
\label{prob:aeg-discriminant}
Under natural completeness and asymptotic hypotheses, determine when the five
loci in Definition~\ref{def:structural-discriminant} admit a locally finite Whitney
stratification and when a Thom-type isotopy lemma applies to the combined AEG data.
\end{problem}
